\documentclass[10pt,a4paper]{amsart}
\usepackage[margin=19mm]{geometry}
\usepackage{amsmath,amssymb,mathtools}
\usepackage[scr=rsfs,cal=euler]{mathalfa}
\usepackage{xfrac}
\usepackage[unicode,hidelinks,bookmarks=false]{hyperref}
\newcommand{\ie}{i.e.\ }
\newcommand{\cf}{cf.\ }
\newcommand{\resp}{respectively\ }
\newcommand{\Subsection}{\subsection}
\providecommand{\nobreakdash}{\nobreak\hbox{-}}
\setlength{\emergencystretch}{2em}
\multlinegap0pt
\usepackage{bm}
\usepackage{bbm}
\usepackage{dsfont}
\usepackage{xspace}
\usepackage{cancel}
\usepackage{mathtools}
\usepackage[shortlabels]{enumitem}
\usepackage{amsthm}
\makeatletter
\@ifundefined{case}{\newtheorem{case}{Case}}{}
\@ifundefined{theorem}{\newtheorem{theorem}{Theorem}}{}
\@ifundefined{corollary}{\newtheorem{corollary}[theorem]{Corollary}}{}
\@ifundefined{proposition}{\newtheorem{proposition}[theorem]{Proposition}}{}
\makeatother
\makeatletter
\@ifundefined{case}{\newenvironment{case}{\refstepcounter{case}\par\medskip\noindent\textbf{Case \arabic{case}. }}{\par\medskip}}{}
\newcommand{\sr}{\hat r}
\newcommand{\sri}{\hat r_{\infty}}
\makeatother
\title[The Erd\H{o}s-Faudree Problems and the Isolate-Free Core]{The Erd\H{o}s-Faudree Problems and the Isolate-Free Core}
\author{Yaping Mao}
\date{Source: arXiv:2604.16012v1 (CC BY 4.0)}
\begin{document}
\maketitle

\noindent\emph{Source and licence.} The Erd\H{o}s-Faudree Problems and the Isolate-Free Core. Version arXiv:2604.16012v1 (CC BY 4.0), distributed under the Creative Commons Attribution 4.0 International licence, as recorded in the arXiv licence metadata for this exact version and shown on the abstract page https://arxiv.org/abs/2604.16012v1. Full rights notice: licenses/arxiv-2604.16012-CC-BY-4.0.txt.\par\smallskip

\noindent\textit{Editorial scope.} Counted unit: the Theorem whose source label is \texttt{thm:connected-classification} in the article (source0.tex lines 506--531, source label \texttt{thm:connected-classification}), delivered together with the article's own complete original proof of it (source0.tex lines 533--618, one \texttt{proof} environment of 86 lines). No mathematics is written, repaired or re-derived: statement and proof are the article's own text line for line. Required context retained uncounted: cor:fekete-reduction (lines 255--271); prop:delta-lower (lines 402--411); eq:basic-params (lines 450--456); prop:bundle-bounds (lines 458--468). Every definition, display and label of those lines is retained unchanged. Omitted from this package and not redistributed: 1--254, 272--401, 412--449, 457--457, 469--505, 532--532, 619--926. Nothing inside a retained excerpt is omitted or altered: every retained body is delivered exactly as the article's lines are written, so no omission receipt of any kind accompanies this package. Citations: the retained excerpts carry no citation command at all, so no bibliography entry of the article is transcribed and no reference is redistributed. Reference adaptation: none was needed and none was made; every reference inside the delivered unit resolves inside it, so no unresolved reference or citation is produced. Printed slips kept literal: no printed word, symbol, hypothesis or number of the counted statement or of its proof was altered. Layout and class: the article's own class is \texttt{article} with packages including inputenc, fontenc, amsmath, amssymb, amsthm, enumitem, cite; the wrapper is the lane's pinned \texttt{amsart} editorial wrapper at 10pt on A4, which restates the theorem-like environments the retained text needs and adds the article's own macro definitions that the retained text needs (\texttt{\textbackslash sr}, \texttt{\textbackslash sri}), with extra wrapper packages (bm, bbm, dsfont, xspace, cancel, mathtools, enumitem). The delivered numbering is the unit's own. Assets: no retained span contains a figure environment, a graphics-inclusion command, a PDF-page inclusion, a table or a TikZ picture; no image file is copied into this package, the package declares no asset dependency and no image file is redistributed with it. Licence: version arXiv:2604.16012v1 of this article is distributed under the Creative Commons Attribution 4.0 International licence, as recorded in the arXiv licence metadata for this exact version and shown on the abstract page https://arxiv.org/abs/2604.16012v1. A full rights notice is delivered at licenses/arxiv-2604.16012-CC-BY-4.0.txt. Characters: every retained excerpt is pure ASCII, so the delivered engine, XeLaTeX with its default Latin Modern text font, reports no missing character.
\begin{corollary}[Fekete reduction]\label{cor:fekete-reduction}
For every nonempty graph $G$,
\[
\sri(G)=\inf_{t\ge 1}\frac{\sr(tK_2,G)}{t|E(G)|}.
\]
Moreover,
$0 \le \sri(G)\le 1$,
and
$\sri(G)=1$
if and only if 
$\sr(tK_2,G)=t|E(G)|$
for every $t\ge 1$.
In particular, for every integer $t\ge 1$,
\[
\sri(G)\le \frac{\sr(tK_2,G)}{t|E(G)|}.
\]
\end{corollary}

\begin{proposition}[A degree lower bound]\label{prop:delta-lower}
For every nonempty graph $G$ and every integer $t\ge 1$,
\[
\sr(tK_2,G)\ge (t-1)\Delta(G)+|E(G)|.
\]
Consequently,
\[
\sri(G)\ge \frac{\Delta(G)}{|E(G)|}.
\]
\end{proposition}

\begin{equation}\label{eq:basic-params}
|V(B)|=n,
\qquad
|E(B)|=k^2+\sum_{i=1}^k\ell_i=n+k(k-2),
\qquad
\Delta(B)=k+\ell_1.
\end{equation}

\begin{proposition}[A two-sided estimate for the bundled bipartite graphs]
\label{prop:bundle-bounds}
Let $B=B(\ell_1,\dots,\ell_k)$ with parameters as above. Then for every integer $t\ge 1$,
\[
(t-1)(k+\ell_1)+n+k(k-2)
\le
\sr(tK_2,B)
\le
n+(t-1)\ell_1+(k+t-1)^2-2k.
\]
\end{proposition}

\begin{theorem}[Connected bipartite limit classification]\label{thm:connected-classification}
For every $\alpha\in[0,1]$ there exists $N_0$ and a family of connected bipartite graphs
$(G_N)_{N\ge N_0}$ such that $|V(G_N)|=N$ and
\[
\sri(G_N)\to \alpha.
\]
More precisely:
\begin{enumerate}[label={\rm (\roman*)}]
\item If $\alpha\in(0,1]$, one may choose a family with
\[
\begin{aligned}
|E(G_N)|&=N+O_{\alpha}(1),\\
\Delta(G_N)&=\alpha N+O_{\alpha}(1),\\
\sri(G_N)&=\alpha+O_{\alpha}(N^{-1/2}).
\end{aligned}
\]
\item If $\alpha=0$ and $t$ is sufficiently large, one may choose a family with
\[
\begin{aligned}
|E(G_N)|&\le 2N,\\
\Delta(G_N)&\le 2\sqrt N+2,\\
\sri(G_N)&=O(N^{-1/2}).
\end{aligned}
\]
\end{enumerate}
\end{theorem}

\begin{proof}
We treat the cases $\alpha>0$ and $\alpha=0$ separately.
\setcounter{case}{0}
\begin{case}
$\alpha\in(0,1]$.    
\end{case}

Choose an integer $k\ge 2$ such that $k\alpha \geq 1$. For each sufficiently large $N$, let
$s_N=N-2k$ and 
$\ell_1=\lfloor \alpha s_N\rfloor$.
Since $k\alpha \geq 1$, we have $s_N-\ell_1\le (k-1)\ell_1$ for all sufficiently large $N$. Hence
one can choose integers
\[
\ell_1\ge \ell_2\ge \cdots\ge \ell_k\ge 0
\qquad\text{with}\qquad
\sum_{i=1}^k\ell_i=s_N.
\]
Set
$G_N=B(\ell_1,\dots,\ell_k)$.
Then $G_N$ is connected and bipartite, $|V(G_N)|=N$, and by \eqref{eq:basic-params},
\[
|E(G_N)|=N+k(k-2)=N+O_{\alpha}(1),
\qquad
\Delta(G_N)=k+\ell_1=\alpha N+O_{\alpha}(1).
\]

By Proposition~\ref{prop:delta-lower}, we have 
\[
\sri(G_N)\ge \frac{k+\ell_1}{N+k(k-2)}=\alpha+O_{\alpha}(N^{-1}).
\]
On the other hand, Corollary~\ref{cor:fekete-reduction} and
Proposition~\ref{prop:bundle-bounds} give, for every integer $t\ge 1$,
\[
\sri(G_N)
\le
\frac{N+(t-1)\ell_1+(k+t-1)^2-2k}{t\bigl(N+k(k-2)\bigr)}.
\]
Since $k$ is fixed and $\ell_1=\alpha N+O_{\alpha}(1)$, the right-hand side is
\[
\alpha+O_{\alpha}\!\left(\frac{1}{t}+\frac{t}{N}\right),
\]
uniformly for all sufficiently large $N$ and all integers $t\ge 1$. Taking the infimum over
all integers $t\ge 1$ and using
\[
\inf_{t\ge 1}\left(\frac{1}{t}+\frac{t}{N}\right)=O(N^{-1/2}),
\]
we obtain
$\sri(G_N)\le \alpha+O_{\alpha}(N^{-1/2})$.
Combining the lower and upper bounds gives $\sri(G_N)\to\alpha$.

\begin{case}
$\alpha=0$.   
\end{case}

For each sufficiently large $N$, let
$q=\lfloor \sqrt N\rfloor$,
$s_N=N-2q$,
and write
$s_N=aq+r
~(0\le r<q)$.
Define
$\ell_1=\cdots=\ell_r=a+1$,
$\ell_{r+1}=\cdots=\ell_q=a$,
and let
$G_N:=B(\ell_1,\dots,\ell_q)$.
Then $G_N$ is connected and bipartite and has $|V(G_N)|=N$. Moreover,
\[
\ell_1\le \left\lceil\frac{N-2q}{q}\right\rceil +1\le \sqrt N+2,
\]
so by \eqref{eq:basic-params},
$|E(G_N)|=N+q(q-2)\le 2N$ and 
$\Delta(G_N)=q+\ell_1\le 2\sqrt N+2$.
Again Corollary~\ref{cor:fekete-reduction} and
Proposition~\ref{prop:bundle-bounds} imply that for every integer $t\ge 1$,
\[
\sri(G_N)
\le
\frac{N+(t-1)\ell_1+(q+t-1)^2-2q}{t\bigl(N+q(q-2)\bigr)}.
\]
Since $q=\lfloor\sqrt N\rfloor$ and $\ell_1\le \sqrt N+2$, the right-hand side is
\[
O\!\left(\frac{1}{2t}+\frac{1}{2\sqrt N}+\frac{t}{2(N-\sqrt N)}\right),
\]
uniformly for all sufficiently large $N$ and $t$.
Hence $\sri(G_N)\to 0$.
\end{proof}

\end{document}
